% Corrected French for printed p.209.
% Four print-level repairs are recorded.
\clearpage
\SourcePage{209}

\LayerTarget{tohoku:conclusion:theorem-5.5.3:p1}
\noindent\textit{précédent est injectif si $n=1$.} De plus :

\LayerTarget{tohoku:corollary:1-theorem-5.5.3}
\textsc{Corollaire 1.} \textit{Si les $U_{i_0\ldots i_p}$ sont tous
$A$-acycliques, alors les homomorphismes \textup{(5.5.5)} sont des isomorphismes.}

\LayerTarget{tohoku:corollary:2-theorem-5.5.3}
\textsc{Corollaire 2.} \textit{Le théorème précédent est encore
valable si ${\mathbf U}$ est un $G$-recouvrement fermé, pourvu que l'on soit dans
l'un ou l'autre cas suivant : a) ${\mathbf U}$ est localement fini et $X$
paracompact, b) ${\mathbf U}$ est fini.}

\LayerTarget{tohoku:proof:corollary-2-theorem-5.5.3}
Dans le cas a), en vertu de ce qui est censé avoir été vu dans 3.8 (voir
remarque de 3.8), les conditions (i) et (ii) du lemme 5.5.1 sont encore
satisfaites (pour (i), c'est d'ailleurs trivial sans hypothèse de
paracompacité, grâce au fait que ${\mathbf U}$ est localement fini), et la
formule $R^q\mathbf K(A)=C({\mathbf U},H^q(A))$ est encore valable. Dans le cas
b), il faut suivre la méthode de Godement (qu'il développe pour $G$ réduit à
l'unité), et que nous allons esquisser, à cause des possibilités d'application
à la Géométrie Algébrique abstraite. Tout d'abord, on introduit, suivant une
idée de P. Cartier, le complexe de faisceaux ${\mathbf C}({\mathbf U},A)$, dont le complexe
des sections sur un ouvert $V$ quelconque est par définition
$C({\mathbf U}|V,A|V)$ (où le signe $|$ indique comme toujours qu'on prend la
restriction à l'ensemble indiqué après ce signe). On prouve que c'est là une
\emph{résolution} de $A$ (voir livre de Godement [9] pour des détails). Nous allons lui
appliquer le corollaire au lemme 5.5.2 (les lettres dans l'énoncé de ce lemme
ayant la même signification que plus haut). Il reste seulement à prouver que
$R^qF(\mathrm C)=H^q(X,{\mathbf C}({\mathbf U},A))$ est isomorphe (fonctoriellement) à
$C({\mathbf U},H^q(A))$. Cela n'offre pas de difficulté, grâce au fait que
${\mathbf U}$ est \emph{fini}, et en utilisant le théorème 3.5.1.

\LayerTarget{tohoku:definition:fixed-point-free-g-cover}
Nous dirons qu'un \emph{$G$-recouvrement} ${\mathbf U}=(U_i)_{i\in I}$ est ``\emph{sans point fixe}'' si $G$ opère dans $I$ sans point fixe, i. e. si $g\ne e$ implique
$g.\, i\ne i$ pour tout $i$. Si alors $P$ est un $G$-préfaisceau quelconque,
alors pour tout $n$, on a

\LayerTarget{tohoku:equation:5.5.6}
\[
\tag{5.5.6}
C^n({\mathbf U},P)=
\prod_{(i_0,\ldots,i_n)}\ \prod_{g\in G}
P(g.\, U_{i_0\ldots i_n}),
\]

\LayerTarget{tohoku:calculation:fixed-point-free-g-cover-cochains}
le premier signe $\displaystyle\prod$ étant étendu à un système complet de représentants
$(i_0,\ldots,i_n)$ de $I^{n+1}/G$. Dans le deuxième produit, tous les termes
sont canoniquement isomorphes à $P(U_{i_0\ldots,i_n})$, de sorte que ce produit
s'identifie au groupe des applications de $G$ dans le groupe fixe
$P(U_{i_0\ldots i_n})$, $G$ y opérant par translations gauches. Il est bien
connu qu'un tel $G$-module est $\Gamma^G$-acyclique, donc il en est de même de
$C^n({\mathbf U},P)$. On fera cependant attention que la formule (5.5.6) n'est
valable qu'en supposant que $C({\mathbf U},P)$ désigne le complexe des cochaines
\emph{quelconques} (non nécessairement alternées) de ${\mathbf U}$ à coéfficients dans
$P$ (alors que les considérations antérieures valaient indifféremment pour les
cochaines quelconques, ou les cochaines alternées). En utilisant la première
suite spectrale d'hyperhomologie du foncteur $\Gamma^G$ par rapport au complexe
$C({\mathbf U},P)$, on trouve alors :

\LayerTarget{tohoku:proposition:5.5.4}
\textsc{Proposition 5.5.4.} \textit{Si ${\mathbf U}$ est un
$G$-recouvrement sans points fixes, alors on a pour tout $G$-préfaisceau $P$ :}

\LayerTarget{tohoku:equation:5.5.7}
\[
\tag{5.5.7}
H^n({\mathbf U},G,P)=H^n\bigl((C({\mathbf U},P))^G\bigr)
\]

\TohokuContinuation{210}
